\section{Discussion}
\label{sec:discussion}

\paragraph{What the evidence supports.} In two models that report both the current and the initial state
correctly, an obsolete state still moves routing decisions toward the action it implies. The net switch
rate is 0.158 (Qwen3-8B) and 0.165 (Qwen3-14B) in the histories where the model retrieved both states and
decided correctly without history. Stale-value interference, documented for value queries
\citep{guo2026context,wang2026unable}, thus also reaches single-step policy decisions in open models after a
single update, on items where retrieval is at ceiling. Retrieval accuracy would not reveal this gap. The
controls separate two sources of the effect. Supersession contributes: in routing, a superseded
assignment of the queried attribute moved preference more than an updated assignment of an unrelated
attribute in both models (1.224 and 0.866 nats). But the same word, mentioned once by the same entity and
never assigned, moves decisions about as much in Qwen3-14B and more in Qwen3-8B. The interference
therefore does not require supersession. It behaves like entrainment by decision-relevant words attached
to the queried entity \citep{niu2025entrainment}, with supersession adding to it. How the two compare is
model-specific.

\paragraph{Consequences for evaluation and mitigation.} Two practical points follow. First, stale-state
benchmarks that vary only update histories cannot tell supersession from incidental mention. A matched
mention condition is cheap and changed our conclusion. Second, mitigations should be evaluated the same
way: with a matched mention control and decision-level outcomes, not retrieval accuracy alone. We did not
test a mitigation. An evaluation should also check historical retrieval, which the model performs well
and which a blunt fix could damage.

\paragraph{Scores and behavior.} Two score-level results caution against reading log-probability shifts
as behavior. An unattached word shifted preferences by 2.5--3.0 nats but changed few decisions (net switch rates
0.039 and 0.062).
And retrieval-score sensitivity to the edit, although large, added nothing beyond the pre-decision
features in predicting which decisions failed. A
preference shift matters behaviorally only when it crosses the item's decision margin, so it can be large
and still inconsequential. Studies that measure only scores, or only accuracy, see different parts of the
same effect.

\paragraph{Alternative explanations.} The constructions differ in syntax around the edited word, so the
contrasts compare whole constructions. Correct direct retrieval also does not show that the current state
was the one used at decision time, so the retrieval checks do not rule out stale binding in the sense of
\citet{guo2026context}. The mention construction could be read as more salient. Its
advantage varied with template in Qwen3-8B and did not appear in Qwen3-14B, which fits a surface
explanation better than a categorical one. The other
entity's current state always implied the wrong action. Confusing the two entities therefore always
yields the wrong action. This may contribute to baseline errors, but it cancels in every paired contrast. Natural state words
carry prior meanings, but the pattern does not depend on them. For Qwen3-8B, in all 12 task $\times$ cell $\times$
vocabulary strata (64 histories each), $D_{\mathrm{superseded}}$ was positive (2.53--4.53 nats).
$\Delta_{\mathrm{mention}}$ was negative in all 12, with intervals excluding zero in 11, including all six with opaque labels
(Table~\ref{tab:strata}).
